%! AP Statistics — Unit 5, Lesson 5.5 — Homework (Answer Key)
\documentclass[10pt]{article}
\usepackage{apstats-article}
\usepackage{apstats-key}

\begin{document}

\pageheader{Unit 5, Lesson 5.5}{Homework --- Answer Key}
\namedateperiod

\vspace{0.15in}

\begin{enumerate}

  \item High school athletes, $p = 0.38$, $n = 180$.
        \begin{enumerate}[label=\alph*.]
          \item \ans{Random: stated. 10\%: $180 \le 10\%$ of high school students nationally --- yes. Large Counts: $np = 180(0.38) = 68.4 \ge 10$; $n(1-p) = 180(0.62) = 111.6 \ge 10$. All conditions met.}
          \item \ans{$\mu_{\hat{p}} = 0.38$; $\sigma_{\hat{p}} = \sqrt{0.38(0.62)/180} \approx 0.03619$.}
          \item \ans{$z = (0.42-0.38)/0.03619 \approx 1.105$; $P(\hat{p}>0.42) \approx 1 - 0.8654 = 0.1346$. About 13.5\%.}
          \item \ans{$z = \text{invNorm}(0.05) \approx -1.645$; $c = 0.38 + (-1.645)(0.03619) \approx 0.320$. There is a 5\% chance that the sample proportion of athletes is less than 32.0\%.}
        \end{enumerate}

  \item Public transit, $p = 0.15$.
        \begin{enumerate}[label=\alph*.]
          \item \ans{$np = 40(0.15) = 6 < 10$. Large Counts condition \emph{not} met. Normal model is NOT appropriate for $n = 40$.}
          \item \ans{$n=100$: $np=15 \ge 10$; $n(1-p)=85 \ge 10$; 10\%: yes. Conditions met. $\sigma_{\hat{p}} = \sqrt{0.15(0.85)/100} \approx 0.03571$. $z_1 = (0.10-0.15)/0.03571 \approx -1.400$; $z_2 = (0.20-0.15)/0.03571 \approx 1.400$; $P \approx 0.9192 - 0.0808 = 0.8384$. About 83.8\%.}
        \end{enumerate}

  \item Red chips, $p = 0.60$.
        \begin{enumerate}[label=\alph*.]
          \item \ans{$n=25$: $\sigma = \sqrt{0.60(0.40)/25} = \sqrt{0.0096} \approx 0.0980$. $n=100$: $\sigma \approx 0.0490$. $n=400$: $\sigma \approx 0.0245$.}
          \item \ansline{$\sigma_{\hat{p}}$ decreases as $n$ increases (halves when $n$ quadruples). Larger samples give more precise estimates of the true proportion.}
        \end{enumerate}

\end{enumerate}

\begin{extensionbox}
\textbf{Extension.}
\ans{$\sigma_{\hat{p}} = \sqrt{0.52(0.48)/600} \approx 0.02040$. $z = (0.50-0.52)/0.02040 \approx -0.980$; $P(\hat{p} < 0.50) \approx 0.1635$. Even with 52\% true support, there is about a 16.4\% chance the poll shows her trailing. This illustrates that polling results can be misleading due to sampling variability alone.}
\end{extensionbox}

\begin{reflectionbox}
\textbf{Preview:} Lesson 5.6 covers sampling distributions for $\hat{p}_1 - \hat{p}_2$.
\end{reflectionbox}

\end{document}
